\documentclass[11pt,a4paper]{article}

\usepackage[a4paper,margin=2.6cm]{geometry}
\usepackage{amsmath,amssymb,mathtools}
\usepackage{bm}
\usepackage{microtype}
\usepackage{xcolor}
\usepackage{hyperref}

\hypersetup{
  colorlinks=true,
  linkcolor=blue!55!black,
  urlcolor=blue!55!black,
  citecolor=blue!55!black
}

\setlength{\parindent}{0pt}
\setlength{\parskip}{0.65em}
\linespread{1.08}
\hyphenpenalty=10000
\exhyphenpenalty=10000
\emergencystretch=3em
\numberwithin{equation}{section}

\newcommand{\R}{\mathbb{R}}
\newcommand{\E}{\mathbb{E}}
\newcommand{\Prob}{\mathbb{P}}
\newcommand{\dd}{\,\mathrm{d}}

\title{\bfseries Understanding Diffusion Through Probability Flux and Multi-Sample Paths}
\author{Math Thought Notes}
\date{June 17, 2026}

\begin{document}
\maketitle

\begin{abstract}
Diffusion appears in molecular motion, chemical transport, cell migration, ecological spread, financial perturbations, heat flow, and contaminant transport. In these settings diffusion is more than the operator \(D\Delta\). It is a modelling language that connects microscopic random paths to macroscopic probability distributions. This note develops diffusion from the viewpoint of multi-sample paths, empirical measures, probability density, Fokker--Planck equations, probability flux, and expected absorption time.
\end{abstract}

\section{From One Random Path to Many Sample Paths}

Consider the one-dimensional diffusion process
\begin{equation}
  \dd X_t=\mu(X_t)\dd t+\sigma(X_t)\dd W_t .
  \label{eq:sde-general}
\end{equation}
Its Euler--Maruyama discretisation is
\begin{equation}
  X_{t+\Delta t}
  =X_t+\mu(X_t)\Delta t+\sigma(X_t)\sqrt{\Delta t}\,\xi,
  \qquad \xi\sim N(0,1).
  \label{eq:euler-maruyama}
\end{equation}
The simplest Brownian diffusion is obtained by taking zero drift and constant diffusion strength:
\begin{equation}
  \dd X_t=\sqrt{2D}\dd W_t,
  \qquad
  X_{t+\Delta t}=X_t+\sqrt{2D\Delta t}\,\xi .
  \label{eq:brownian-discrete}
\end{equation}

A single path \(t\mapsto X_t\) describes one possible microscopic history. The object that becomes stable in applications is usually not a single path but the collective behaviour of many independent paths. Let
\begin{equation}
  X_t^{(1)},X_t^{(2)},\ldots,X_t^{(N)}
  \label{eq:sample-paths}
\end{equation}
be \(N\) independent copies of the diffusion in \eqref{eq:sde-general}. Their empirical measure at time \(t\) is
\begin{equation}
  \mu_t^N
  =\frac{1}{N}\sum_{m=1}^{N}\delta_{X_t^{(m)}} .
  \label{eq:empirical-measure}
\end{equation}
For any smooth test function \(\varphi\), the empirical measure gives
\begin{equation}
  \int_{\R}\varphi(x)\dd\mu_t^N(x)
  =\frac{1}{N}\sum_{m=1}^{N}\varphi(X_t^{(m)}).
  \label{eq:empirical-test}
\end{equation}
By the law of large numbers,
\begin{equation}
  \frac{1}{N}\sum_{m=1}^{N}\varphi(X_t^{(m)})
  \longrightarrow
  \E[\varphi(X_t)]
  \qquad \text{as } N\to\infty .
  \label{eq:law-large-numbers}
\end{equation}
If \(X_t\) has density \(p(x,t)\), then
\begin{equation}
  \E[\varphi(X_t)]
  =\int_{\R}\varphi(x)p(x,t)\dd x .
  \label{eq:density-weak-form}
\end{equation}
Thus the large-sample limit of many paths is naturally a probability density.

The same idea appears in a histogram. If the interval \([x_i,x_i+h)\) contains
\begin{equation}
  H_i(t)=\#\{m:X_t^{(m)}\in[x_i,x_i+h)\}
  \label{eq:hist-count}
\end{equation}
sample particles, then
\begin{equation}
  \frac{H_i(t)}{Nh}\approx p(x_i,t).
  \label{eq:hist-density}
\end{equation}
The density \(p(x,t)\) can therefore be interpreted as the statistical shape of many sample paths at a fixed time. Diffusion at the macroscopic level is the collective structure produced by microscopic random paths.

\section{Generator and the Fokker--Planck Equation}

The infinitesimal generator of \eqref{eq:sde-general} is
\begin{equation}
  \mathcal{L}\varphi(x)
  =\mu(x)\varphi'(x)
  +\frac{1}{2}\sigma^2(x)\varphi''(x).
  \label{eq:generator}
\end{equation}
Ito's formula gives
\begin{equation}
  \dd\varphi(X_t)
  =\mathcal{L}\varphi(X_t)\dd t
  +\sigma(X_t)\varphi'(X_t)\dd W_t.
  \label{eq:ito}
\end{equation}
After taking expectation, the martingale term vanishes:
\begin{equation}
  \frac{\dd}{\dd t}\E[\varphi(X_t)]
  =\E[\mathcal{L}\varphi(X_t)].
  \label{eq:expect-generator}
\end{equation}
Using the density \(p(x,t)\), equation \eqref{eq:expect-generator} becomes
\begin{equation}
  \frac{\dd}{\dd t}\int_{\R}\varphi(x)p(x,t)\dd x
  =
  \int_{\R}\mathcal{L}\varphi(x)p(x,t)\dd x.
  \label{eq:weak-forward}
\end{equation}
Integration by parts transfers derivatives from \(\varphi\) to \(p\):
\begin{equation}
  \int_{\R}\mathcal{L}\varphi(x)p(x,t)\dd x
  =
  \int_{\R}\varphi(x)\mathcal{L}^*p(x,t)\dd x,
  \label{eq:adjoint-definition}
\end{equation}
where
\begin{equation}
  \mathcal{L}^*p
  =
  -\frac{\partial}{\partial x}\bigl(\mu(x)p\bigr)
  +\frac{1}{2}\frac{\partial^2}{\partial x^2}
  \bigl(\sigma^2(x)p\bigr).
  \label{eq:adjoint}
\end{equation}
Since \(\varphi\) is arbitrary, the density satisfies the Fokker--Planck equation
\begin{equation}
  \frac{\partial p}{\partial t}
  =
  -\frac{\partial}{\partial x}\bigl(\mu(x)p\bigr)
  +\frac{1}{2}\frac{\partial^2}{\partial x^2}
  \bigl(\sigma^2(x)p\bigr).
  \label{eq:fokker-planck}
\end{equation}
This derivation shows that the Fokker--Planck equation is the distribution-level representation of the SDE. The SDE describes sample paths, while \eqref{eq:fokker-planck} describes the statistical limit of many sample paths.

For Brownian diffusion,
\begin{equation}
  \mu(x)=0,
  \qquad
  \sigma^2(x)=2D.
  \label{eq:brownian-coefficients}
\end{equation}
Then \eqref{eq:fokker-planck} reduces to
\begin{equation}
  \frac{\partial p}{\partial t}
  =D\frac{\partial^2p}{\partial x^2}.
  \label{eq:heat-1d}
\end{equation}
In multiple dimensions this becomes
\begin{equation}
  \frac{\partial p}{\partial t}=D\Delta p.
  \label{eq:heat-nd}
\end{equation}
The heat equation is therefore also the probability-density equation of Brownian motion.

\section{Fundamental Solution and the Spread of a Sample Cloud}

If all particles start at the origin,
\begin{equation}
  X_0=0,
  \label{eq:origin-initial}
\end{equation}
then Brownian diffusion satisfies
\begin{equation}
  X_t=\sqrt{2D}W_t.
  \label{eq:brownian-solution}
\end{equation}
Consequently,
\begin{equation}
  X_t\sim N(0,2Dt),
  \label{eq:brownian-distribution}
\end{equation}
and the density is
\begin{equation}
  p(x,t)
  =
  \frac{1}{\sqrt{4\pi Dt}}
  \exp\left(-\frac{x^2}{4Dt}\right).
  \label{eq:gaussian-kernel}
\end{equation}
This Gaussian kernel is the fundamental solution of
\begin{equation}
  \frac{\partial p}{\partial t}
  =D\frac{\partial^2p}{\partial x^2}
  \label{eq:heat-fundamental}
\end{equation}
with initial condition
\begin{equation}
  p(x,0)=\delta_0(x).
  \label{eq:delta-initial}
\end{equation}
Its moments are
\begin{equation}
  \E[X_t]=0,
  \qquad
  \operatorname{Var}(X_t)=2Dt.
  \label{eq:brownian-moments}
\end{equation}
Thus a cloud of particles released from a single point keeps its centre at the origin while its width grows like
\begin{equation}
  \sqrt{\operatorname{Var}(X_t)}
  =\sqrt{2Dt}.
  \label{eq:diffusive-scale}
\end{equation}
Diffusion is not primarily the movement of the mean position; it is the spreading of a distribution. This interpretation is useful for molecular diffusion, contaminant transport, cell population migration, heat diffusion, and random perturbation propagation.

For a general initial density \(p_0(x)\), the solution is the convolution
\begin{equation}
  p(x,t)
  =
  \int_{\R}
  \frac{1}{\sqrt{4\pi Dt}}
  \exp\left(-\frac{(x-y)^2}{4Dt}\right)
  p_0(y)\dd y.
  \label{eq:heat-convolution}
\end{equation}
This formula expresses diffusion as Gaussian smoothing of the initial distribution.

\section{Probability Flux and Conservation Form}

The Fokker--Planck equation can be written as a conservation law:
\begin{equation}
  \frac{\partial p}{\partial t}
  +\frac{\partial J}{\partial x}=0.
  \label{eq:conservation-form}
\end{equation}
For the one-dimensional diffusion \eqref{eq:sde-general}, the flux is
\begin{equation}
  J(x,t)
  =
  \mu(x)p(x,t)
  -\frac{1}{2}\frac{\partial}{\partial x}
  \bigl(\sigma^2(x)p(x,t)\bigr).
  \label{eq:general-flux}
\end{equation}
For pure diffusion,
\begin{equation}
  J=-D\frac{\partial p}{\partial x}.
  \label{eq:ficks-law}
\end{equation}
This is Fick's law: probability mass flows from regions of high density to regions of low density.

With constant drift \(v\),
\begin{equation}
  \dd X_t=v\dd t+\sqrt{2D}\dd W_t,
  \label{eq:constant-drift-sde}
\end{equation}
the density satisfies
\begin{equation}
  \frac{\partial p}{\partial t}
  =
  -v\frac{\partial p}{\partial x}
  +D\frac{\partial^2p}{\partial x^2},
  \label{eq:advection-diffusion-1d}
\end{equation}
with flux
\begin{equation}
  J=vp-D\frac{\partial p}{\partial x}.
  \label{eq:advection-diffusion-flux}
\end{equation}
The term \(vp\) is advection flux, while \(-D\partial_xp\) is diffusion flux. This form appears naturally in contaminant transport in a flow, cell movement under directional bias, and particles transported by an external velocity field.

On an interval \([a,b]\), the total probability mass is
\begin{equation}
  M(t)=\int_a^b p(x,t)\dd x.
  \label{eq:interval-mass}
\end{equation}
Integrating \eqref{eq:conservation-form} gives
\begin{equation}
  \frac{\dd M}{\dd t}=J(a,t)-J(b,t).
  \label{eq:mass-balance-interval}
\end{equation}
The change of probability mass inside a region is determined entirely by the flux through its boundary.

\section{Reflecting, Absorbing, and Partially Absorbing Boundaries}

Boundary conditions have direct probabilistic meanings. A reflecting boundary has no probability flux through it. On \([0,L]\), reflecting boundaries are written as
\begin{equation}
  J(0,t)=J(L,t)=0.
  \label{eq:reflecting-flux}
\end{equation}
For pure diffusion this is equivalent to the Neumann condition
\begin{equation}
  \frac{\partial p}{\partial x}(0,t)=0,
  \qquad
  \frac{\partial p}{\partial x}(L,t)=0.
  \label{eq:reflecting-neumann}
\end{equation}
Therefore
\begin{equation}
  \frac{\dd}{\dd t}\int_0^L p(x,t)\dd x=0.
  \label{eq:reflecting-mass}
\end{equation}
Total probability is conserved. Reflecting boundaries model impermeable walls, closed containers, or symmetry surfaces with zero net flux.

At an absorbing boundary, a particle leaves the system after reaching the boundary. A typical absorbing condition is
\begin{equation}
  p(x,t)=0,
  \qquad x\in\partial\Omega_{\mathrm{abs}}.
  \label{eq:absorbing-boundary}
\end{equation}
Then
\begin{equation}
  \frac{\dd}{\dd t}\int_{\Omega}p(x,t)\dd x\le 0.
  \label{eq:absorbing-mass-decay}
\end{equation}
Let
\begin{equation}
  \tau=\inf\{t:X_t\in\partial\Omega_{\mathrm{abs}}\}.
  \label{eq:hitting-time}
\end{equation}
The probability mass remaining inside \(\Omega\) is the survival probability:
\begin{equation}
  \Prob(\tau>t)=\int_{\Omega}p(x,t)\dd x.
  \label{eq:survival-probability}
\end{equation}
The lost mass is
\begin{equation}
  1-\int_{\Omega}p(x,t)\dd x=\Prob(\tau\le t).
  \label{eq:absorbed-probability}
\end{equation}
Thus an absorbing-boundary density is naturally a sub-probability density.

A partially absorbing boundary lies between reflection and absorption. A common Robin condition is
\begin{equation}
  -D\frac{\partial p}{\partial n}
  =\kappa p,
  \qquad x\in\partial\Omega,
  \label{eq:robin-boundary}
\end{equation}
where \(\kappa\ge0\) is the boundary absorption strength. The case \(\kappa=0\) gives reflection; the limit \(\kappa\to\infty\) approaches full absorption. Such boundary conditions occur in surface reactions, receptor binding, membrane absorption, and catalytic interfaces.

\section{Expected Absorption Time}

For absorption problems, an important quantity is the expected absorption time
\begin{equation}
  T(x)=\E_x[\tau],
  \label{eq:mean-absorption-time}
\end{equation}
where the process starts at \(X_0=x\). For generator \(\mathcal{L}\), the mean absorption time satisfies the backward equation
\begin{equation}
  \mathcal{L}T(x)=-1,
  \qquad x\in\Omega,
  \label{eq:backward-mean-time}
\end{equation}
with absorbing boundary condition
\begin{equation}
  T(x)=0,
  \qquad x\in\partial\Omega_{\mathrm{abs}}.
  \label{eq:mean-time-absorbing-bc}
\end{equation}
If part of the boundary is reflecting, then
\begin{equation}
  \frac{\partial T}{\partial n}=0,
  \qquad x\in\partial\Omega_{\mathrm{ref}}.
  \label{eq:mean-time-reflecting-bc}
\end{equation}
For Brownian diffusion,
\begin{equation}
  \mathcal{L}=D\Delta,
  \label{eq:brownian-generator}
\end{equation}
so
\begin{equation}
  D\Delta T(x)=-1.
  \label{eq:brownian-mean-time}
\end{equation}

On the interval \([0,L]\) with both endpoints absorbing,
\begin{equation}
  DT''(x)=-1,
  \qquad
  T(0)=T(L)=0.
  \label{eq:two-absorbing-bvp}
\end{equation}
Integrating twice gives
\begin{align}
  T''(x)&=-\frac{1}{D}, \label{eq:two-absorbing-step1}\\
  T'(x)&=-\frac{x}{D}+C_1, \label{eq:two-absorbing-step2}\\
  T(x)&=-\frac{x^2}{2D}+C_1x+C_2. \label{eq:two-absorbing-step3}
\end{align}
The condition \(T(0)=0\) gives
\begin{equation}
  C_2=0.
  \label{eq:c2-zero}
\end{equation}
The condition \(T(L)=0\) gives
\begin{equation}
  -\frac{L^2}{2D}+C_1L=0,
  \qquad
  C_1=\frac{L}{2D}.
  \label{eq:c1-two-absorbing}
\end{equation}
Therefore
\begin{equation}
  T(x)=\frac{x(L-x)}{2D}.
  \label{eq:two-absorbing-solution}
\end{equation}
The maximum occurs at the midpoint:
\begin{equation}
  T\left(\frac{L}{2}\right)
  =\frac{L^2}{8D}.
  \label{eq:two-absorbing-midpoint}
\end{equation}

If the left endpoint is reflecting and the right endpoint is absorbing, then
\begin{equation}
  T'(0)=0,
  \qquad
  T(L)=0,
  \qquad
  DT''(x)=-1.
  \label{eq:reflect-absorb-bvp}
\end{equation}
The first integration gives
\begin{equation}
  T'(x)=-\frac{x}{D}+C_1.
  \label{eq:reflect-absorb-step1}
\end{equation}
Using \(T'(0)=0\) gives
\begin{equation}
  C_1=0.
  \label{eq:reflect-absorb-c1}
\end{equation}
Using \(T(L)=0\) then gives
\begin{equation}
  T(x)=\frac{L^2-x^2}{2D}.
  \label{eq:reflect-absorb-solution}
\end{equation}
In particular,
\begin{equation}
  T(0)=\frac{L^2}{2D}.
  \label{eq:reflect-absorb-origin}
\end{equation}
Reflection increases the expected absorption time because the particle cannot leave through the reflecting endpoint.

\section{Applied Interpretation}

In molecular diffusion, \(p(x,t)\) can represent the probability density of a single molecule or, after scaling by the number of molecules, a concentration profile. If a boundary is absorbing, the same formalism describes molecules captured by a surface, bound by receptors, absorbed by a membrane, or entering a reactive region. The boundary flux gives an absorption rate.

In ecology and cell migration, diffusion models random movement, while drift represents chemotaxis, environmental gradients, or external flow. The advection--diffusion equation
\begin{equation}
  \frac{\partial p}{\partial t}
  =
  -\nabla\cdot(vp)+D\Delta p
  \label{eq:advection-diffusion-nd}
\end{equation}
has flux
\begin{equation}
  J=vp-D\nabla p.
  \label{eq:advection-diffusion-nd-flux}
\end{equation}
The term \(vp\) is directed transport, while \(-D\nabla p\) is random spreading.

For a reaction--diffusion model
\begin{equation}
  \frac{\partial u}{\partial t}=D\Delta u+f(u),
  \label{eq:reaction-diffusion}
\end{equation}
the term \(D\Delta u\) represents spatial spreading and \(f(u)\) represents local growth, decay, reaction, or interaction. If \(u\) is a probability density, then
\begin{equation}
  \int_{\Omega}u(x,t)\dd x
  \label{eq:total-u}
\end{equation}
is total probability mass. If \(u\) is a concentration, the same integral is total amount.

For the linear decay model
\begin{equation}
  \frac{\partial u}{\partial t}=D\Delta u-ku,
  \label{eq:diffusion-decay}
\end{equation}
with no-flux boundary, the total mass satisfies
\begin{equation}
  \frac{\dd}{\dd t}\int_{\Omega}u(x,t)\dd x
  =
  -k\int_{\Omega}u(x,t)\dd x.
  \label{eq:decay-mass-balance}
\end{equation}
If absorption at the boundary is also present, then
\begin{equation}
  \frac{\dd}{\dd t}\int_{\Omega}u(x,t)\dd x
  =
  -\int_{\partial\Omega}J\cdot n\dd S
  -k\int_{\Omega}u(x,t)\dd x.
  \label{eq:boundary-reaction-mass-balance}
\end{equation}
This mass balance combines diffusion, boundary absorption, and reaction decay in one expression.

\section{Mathematical Structure}

Diffusion has three linked levels. The first is the path level:
\begin{equation}
  \dd X_t=\mu(X_t)\dd t+\sigma(X_t)\dd W_t.
  \label{eq:summary-path}
\end{equation}
The second is the distribution level:
\begin{equation}
  \frac{\partial p}{\partial t}=\mathcal{L}^*p.
  \label{eq:summary-distribution}
\end{equation}
The third is the event level, represented here by hitting time:
\begin{equation}
  T(x)=\E_x[\tau],
  \qquad
  \mathcal{L}T=-1.
  \label{eq:summary-event}
\end{equation}
Multi-sample paths connect the path level to the distribution level. Probability flux connects density evolution to boundary behaviour. Expected absorption time connects random paths to boundary events.

The objects can be placed in one framework:
\begin{equation}
  X_t
  \quad\longleftrightarrow\quad
  \mu_t^N
  \quad\longleftrightarrow\quad
  p(x,t)
  \quad\longleftrightarrow\quad
  J(x,t)
  \quad\longleftrightarrow\quad
  T(x).
  \label{eq:object-chain}
\end{equation}
The path gives microscopic random motion, the empirical measure gives multi-sample statistics, the density gives probability distribution, the flux gives probability movement, and the expected absorption time gives the average time scale of a boundary event. Diffusion appears so widely because these objects can be connected within the same mathematical structure.

\end{document}
